\subsection{Exact closure of discounted first moments}\label{subsec: momentClosure}
This section gives a deterministic formulation of the four-process
representation. The closure is exact for a prescribed, possibly nonconstant,
local volatility. A converse below identifies sufficiently regular solutions
of the nonlinear closed equations with calibrated stochastic models. It does
not assert existence of such nonlinear solutions.

\subsubsection{Coordinates and operators}
Use the domestic risk-neutral measure and set
\[
Z_t=(X_t,U_t,V_t)=(\log S_t,x_t^d,x_t^f),\qquad
s(t,x)=\sigma(t,e^x),\qquad g(t,x)=\partial_xs(t,x).
\]
Here $X$ denotes log spot, not the three-vector used in the manuscript's
variational construction. Write $\nu_i=\sigma_r^i$, $r_d=\eta_d(t)+u$,
$r_f=\eta_f(t)+v$, where the curve-fitting shifts are
\[
\eta_i(t)=r_i^0(t)+\tfrac12\nu_i^2 B_{\kappa_i}(t)^2,
\qquad B_k(t)=\int_0^t e^{-ka}da.
\]
In particular $B_0(t)=t$. Put
\[
\beta=\rho_{Sd}\nu_d,\quad \gamma=\rho_{Sf}\nu_f,\quad
c=\rho_{df}\nu_d\nu_f,\quad k=\kappa_d+\kappa_f,\quad
v_i(t)=\nu_i^2 B_{2\kappa_i}(t).
\]
The drift, covariance matrix and covariance with the spot Brownian motion are
\begin{equation}\label{eq: closureState}
b=\begin{pmatrix}r_d-r_f-\tfrac12s^2\\-\kappa_du\\-\kappa_fv-\gamma s\end{pmatrix},
\quad Q=\begin{pmatrix}s^2&\beta s&\gamma s\\\beta s&\nu_d^2&c\\\gamma s&c&\nu_f^2\end{pmatrix},
\quad h=\begin{pmatrix}s\\\beta\\\gamma\end{pmatrix}.
\end{equation}
Thus $d\langle Z,W^S\rangle_t=h(t,Z_t)dt$. Define the backward killed operator
and its distributional adjoint by
\begin{align}
\mathcal K_s\varphi&=b\cdot\nabla\varphi+
\tfrac12\sum_{a,b=1}^3Q_{ab}\partial_{ab}\varphi-r_d\varphi,
\label{eq: closureOperator}\\
\mathcal K_s^*w&=-\nabla\cdot(bw)+
\tfrac12\sum_{a,b=1}^3\partial_{ab}(Q_{ab}w)-r_dw.\nonumber
\end{align}
The off-diagonal second derivatives each occur twice inside the sum.

Order the auxiliary variables as $Y=(\ell^d,\zeta^f,\ell^f,\zeta^d)^\top$.
Their equations take the affine form
\begin{equation}\label{eq: closureAffine}
dY=(AY+f)dt+GY\,dW^S,
\end{equation}
where
\[
A=\begin{pmatrix}
-\kappa_d-sg&-1&0&0\\
-\gamma g&-k&0&0\\
0&0&-\kappa_f&1\\
0&0&\beta g&-k
\end{pmatrix},\qquad
f=\begin{pmatrix}v_d+\beta s\\c\\-v_f+\gamma s\\c\end{pmatrix},
\qquad G=\operatorname{diag}(g,0,g,0).
\]
The initial conditions are $Z_0=z_0=(\log S_0,0,0)$ and $Y_0=0$.

\subsubsection{The closed forward equations}
For $D_t=\exp(-\int_0^t r_d(a,Z_a)da)$ define a positive measure $p_t$
and four signed measures $m_{i,t}$ on $\mathbb R^3$ by
\begin{equation}\label{eq: closureMeasures}
\langle\varphi,p_t\rangle=\mathbb E^d[D_t\varphi(Z_t)],\qquad
\langle\varphi,m_{i,t}\rangle=\mathbb E^d[D_tY_{i,t}\varphi(Z_t)].
\end{equation}
These are discounted measures, not probability densities; $m_i$ need not be
positive. The same letters denote densities when they exist.

\begin{thm}[Exact discounted first-moment closure]\label{thm: momentClosure}
Suppose the state and affine auxiliary equations have nonexplosive solutions.
Assume $\mathbb E[D_t(1+|Y_t|)]<\infty$ and, for every compactly supported
smooth test function, integrability of the drift terms and the true-martingale
property of the local martingales in It\^o's formulas for
$D\varphi(Z)$ and $DY_i\varphi(Z)$. Then, in the distributional sense,
\begin{align}
\partial_t p&=\mathcal K_s^*p,\label{eq: closureFive}\\
\partial_t m_1&=\mathcal K_s^*m_1-(\kappa_d+sg)m_1-m_2
 +(v_d+\beta s)p-\nabla\cdot(ghm_1),\nonumber\\
\partial_t m_2&=\mathcal K_s^*m_2-km_2-\gamma gm_1+cp,\nonumber\\
\partial_t m_3&=\mathcal K_s^*m_3-\kappa_fm_3+m_4
 +(-v_f+\gamma s)p-\nabla\cdot(ghm_3),\nonumber\\
\partial_t m_4&=\mathcal K_s^*m_4-km_4+\beta gm_3+cp,\nonumber
\end{align}
with $p_0=\delta_{z_0}$ and $m_{i,0}=0$.
No second or higher moments of $Y$ occur. This assertion does not require
invertibility of the driver correlation matrix or existence of densities.
\end{thm}
\begin{proof}
The finite-variation differential $dD=-r_dDdt$ contributes the killing term.
The product rule for $DY_i\varphi(Z)$ gives the drift
\[
D\left\{Y_i\mathcal K_s\varphi+(AY+f)_i\varphi
            +(GY)_i h\cdot\nabla\varphi\right\}dt.
\]
The last summand is the quadratic covariation of $Y_i$ and $\varphi(Z)$.
Taking expectations gives
\[
\frac{d}{dt}\langle\varphi,m_i\rangle
=\langle\mathcal K_s\varphi,m_i\rangle+
\sum_j\langle A_{ij}\varphi,m_j\rangle+\langle f_i\varphi,p\rangle
+\mathbf1_{i\in\{1,3\}}\langle gh\cdot\nabla\varphi,m_i\rangle.
\]
Distributional integration by parts proves the assertion. Equivalently, the
joint generator preserves functions affine in $Y$. Their second derivatives
with respect to $Y$ vanish, so the quadratic auxiliary diffusion cannot produce
second auxiliary moments. All mixed state--auxiliary derivatives are retained
in the divergence term.
\end{proof}

\paragraph{Explicit shared-noise term.}
For either $m=m_1$ or $m=m_3$,
\[
-\nabla\cdot(ghm)=-\partial_x(sgm)-\beta\partial_u(gm)
                                      -\gamma\partial_v(gm).
\]
Omitting this term is incorrect even when the physical state diffusion is
otherwise implemented exactly. For fixed $s$ the five equations are linear:
one density equation and two pairs of coupled moment equations.

\subsubsection{Uniqueness for a prescribed nonflat coefficient}
The preceding theorem uses only integrability assumptions. The following
stronger, concrete assumptions give a convenient uniqueness class; they are
not claimed to be minimal.

\begin{prop}[Fixed-coefficient uniqueness]\label{prop: closureIdentification}
Fix a finite horizon. Assume the two rate volatilities are strictly positive,
the constant correlation matrix is positive definite, $\kappa_i\geq0$, and
the deterministic shifts are $C^1$.
Suppose $s$ is $C^{1,4}$, is bounded above and away from zero, and its spatial
derivatives through order four are bounded on the horizon. Then the stochastic
measures in (\ref{eq: closureMeasures}) are the unique weakly continuous
measure solution of (\ref{eq: closureFive}) with its specified initial data in
the class
\begin{equation}\label{eq: closureWeightClass}
\sup_{0\leq t\leq T}\int e^{\lambda|z|}
       \left(|p_t|+\sum_{i=1}^4|m_{i,t}|\right)(dz)<\infty
\quad\hbox{for every }\lambda>0.
\end{equation}
Here absolute values denote total variation. In particular the fixed-coefficient
system does not introduce additional solutions in this class. The volatility
may have genuine spot dependence.
\end{prop}
\begin{proof}
The state coefficients have linear growth and bounded spatial derivatives of
positive order needed below. The state SDE has a unique strong solution. The
rate states are Gaussian convolutions plus bounded drift convolutions; log
spot is an integral of these states plus a bounded-volatility martingale and
a bounded drift correction. Hence $Z$ has all exponential moments of $|Z|$,
uniformly on the finite horizon. The matrices $A,G$ and the source $f$ are
bounded. Standard It\^o moment estimates for their linear SDE give all finite
polynomial moments of $Y$. The domestic discount factor and its inverse have
all finite moments because the domestic integrated rate is Gaussian.
H\"older's inequality proves (\ref{eq: closureWeightClass}) for the stochastic
measures and the integrability required in Theorem \ref{thm: momentClosure}.

For uniqueness take arbitrary smooth compactly supported terminal functions
$\psi_0,\ldots,\psi_4$ at time $t$. Starting from $(Z_a,Y_a)=(z,y)$, define
\[
F(a,z,y)=\mathbb E_{a,z,y}\left[
 e^{-\int_a^t r_d(b,Z_b)db}
 \left(\psi_0(Z_t)+\sum_iY_{i,t}\psi_i(Z_t)\right)\right].
\]
The affine auxiliary equation implies
$F(a,z,y)=F_0(a,z)+\sum_i y_i F_i(a,z)$.
Differentiation of the stochastic flow up to order two is legitimate under the
bounded derivative assumptions; linear auxiliary moment estimates also apply
to these variations. It gives spatial derivatives of $F_j$ up to order two
with at most exponential growth in $|z|$. The Markov property and It\^o's
formula give the backward equations
\begin{align*}
\partial_a F_0+\mathcal K_s F_0+\sum_i f_iF_i&=0,\\
\partial_a F_j+\mathcal K_s F_j+\sum_i A_{ij}F_i
       +\mathbf1_{j\in\{1,3\}}gh\cdot\nabla F_j&=0.
\end{align*}
Here $f$ is the source in (\ref{eq: closureAffine}); the terminal conditions
are $F_j(t,\cdot)=\psi_j$.
Time regularity follows from these equations (or from the uniformly elliptic
backward equation). Pair these functions with the difference of two forward
solutions. Their time derivative is zero by duality. To use these noncompact
test functions, multiply by smooth cutoffs tending to one. All cutoff errors
vanish by (\ref{eq: closureWeightClass}), the exponential derivative bounds,
and the at-most-linear state drift and killing term. The initial difference
is zero, so the terminal pairing is zero for every choice of $\psi_j$.
This proves uniqueness of all five measures.
\end{proof}

\subsubsection{Compatibility identities and the spot marginal}
The next identities additionally use the fact that $Y$ consists of the
integrated Malliavin sensitivities, rather than arbitrary affine auxiliaries.
Impose the own-forward-measure duality and integrability assumptions of that
representation; the bounded smooth setting of Proposition
\ref{prop: closureIdentification} is a sufficient setting. Define
$c_d=r_d-r_d^0(t)$ and $c_f=r_f-r_f^0(t)$.

\begin{prop}[Full-state and marginal compatibility]\label{prop: closureCompatibility}
For the stochastic measures, the following are distributional identities:
\begin{align}
c_dp&=-\partial_xm_1-v_d\partial_up-\partial_vm_2,
\label{eq: closureConstraints}\\
c_fp&=-\partial_xm_3-m_3-\partial_um_4-v_f\partial_vp.\nonumber
\end{align}
Let bars denote integration over $(u,v)$, put $n_i=\overline{m_i}$,
$n=n_1-n_3$, and $J_i=\overline{c_i p}$. Then
\begin{equation}\label{eq: closureMarginalConstraints}
J_d=-\partial_xn_1,\qquad J_f=-\partial_xn_3-n_3.
\end{equation}
Consequently the discounted log-spot marginal obeys
\begin{equation}\label{eq: closureMarginal}
\partial_t\bar p=\mathcal B_t\bar p+
 \tfrac12(\partial_{xx}+\partial_x)(s^2\bar p+2n),
\qquad
\mathcal B_t w=-(r_d^0-r_f^0)\partial_xw-r_d^0 w.
\end{equation}
\end{prop}
\begin{proof}
For a smooth test $\varphi$ of all three coordinates, domestic own-forward
duality gives
\[
\mathbb E^d[D c_d\varphi(Z)]
=\mathbb E^d[D(\ell^d\partial_x\varphi+
                         v_d\partial_u\varphi+\zeta^f\partial_v\varphi)].
\]
Foreign own-forward duality gives the analogous identity with tangent vector
$(\ell^f,\zeta^d,v_f)$ and, on changing to the domestic measure, a factor
$S=e^x$ on both sides. Thus
\[
e^xc_fp=-\partial_x(e^xm_3)-\partial_u(e^xm_4)
                         -v_f\partial_v(e^xp).
\]
This proves (\ref{eq: closureConstraints}), including the essential $-m_3$
term. Rate integration, justified by cutoffs and the moment bounds, proves
(\ref{eq: closureMarginalConstraints}). Direct integration of the density
equation gives
\[
\partial_t\bar p=
-\partial_x\left[(r_d^0-r_f^0-\tfrac12s^2)\bar p+J_d-J_f\right]
+\tfrac12\partial_{xx}(s^2\bar p)-r_d^0\bar p-J_d.
\]
Substitution of the two marginal constraints yields
(\ref{eq: closureMarginal}).
\end{proof}

If the respective upper-boundary fluxes vanish, the same identities imply
\begin{equation}\label{eq: closureTailRecovery}
n_1(x)=\int_x^\infty J_d(y)dy,\qquad
n_3(x)=e^{-x}\int_x^\infty e^yJ_f(y)dy.
\end{equation}
Thus the correction can also be recovered from the physical discounted
density alone, using the conventional tail integrals. The moment fields give
an alternative evolution of those quantities; they do not add independent
terminal-law information. The five-field formulation is not a claim that
three-dimensional density-based calibration was previously unavailable.

\subsubsection{Nonlinear calibration and a known-denominator formulation}
Let $q(t,x)=\sigma_{Dup}(t,e^x)^2$ be the prescribed market Dupire variance.
With $\bar p>0$, the exact nonlinear feedback is
\begin{equation}\label{eq: closureFeedback}
s^2(t,x)=q(t,x)-2\frac{n(t,x)}{\bar p(t,x)},\qquad
n(t,x)=\int_{\mathbb R^2}(m_1-m_3)(t,x,u,v)du\,dv.
\end{equation}
There is no log-coordinate Jacobian in this ratio: it cancels between the
two marginals. Both are discounted under the domestic measure.
Equations (\ref{eq: closureFive}) and (\ref{eq: closureFeedback}), with
$g=\partial_xs$, are a closed nonlinear system of five fields on three spatial
coordinates. They are not five equations on log spot alone.

Let $p_*(t,x)$ denote the market discounted log-spot density. For a sufficiently
regular market call surface $C$,
\begin{equation}\label{eq: closureMarketDensity}
p_*(t,x)=e^x\partial_{KK}C(t,e^x),\quad
\partial_t p_* =\mathcal B_t p_*+
              \tfrac12(\partial_{xx}+\partial_x)(q p_*),\quad
p_*(0)=\delta_{\log S_0}.
\end{equation}
Assume it is positive for $t>0$ and that the stated one-dimensional forward
equations have uniqueness in the integrability class under consideration.
For smooth bounded uniformly positive coefficients and deterministic bounded
rates, this is the usual diffusion-law uniqueness class; exponential moments
also justify recovering call prices from the densities.

\begin{thm}[Calibration equivalence]
\label{thm: closureNonlinear}
Suppose the five fields solve (\ref{eq: closureFive}) in the class
(\ref{eq: closureWeightClass}), with densities for positive times and the
specified initial measure data. Suppose the induced deterministic coefficient $s$
satisfies Proposition \ref{prop: closureIdentification}, and the market data
satisfy (\ref{eq: closureMarketDensity}) and the uniqueness assumptions above.
Either of the feedback relations
\begin{align}
s^2&=q-2n/\bar p,\label{eq: closureTwoFeedbacks}\\
\text{or}\qquad s^2&=q-2n/p_*\label{eq: closureKnownDenominator}
\end{align}
for $t>0$ implies $\bar p=p_*$ and exact calibration of the associated
three-factor model. In this regularity class the two formulations are
equivalent, and their solutions are precisely the stochastic moment fields
of calibrated coefficients. This is an equivalence conditional on a solution,
not an existence or uniqueness assertion for the nonlinear system.
\end{thm}
\begin{proof}
Once the candidate $s$ is fixed, Proposition \ref{prop: closureIdentification}
identifies all five fields with the stochastic measures of the corresponding
strong state and auxiliary solution. Proposition \ref{prop: closureCompatibility}
therefore applies; it is not an additional freely imposed constraint.
With the first feedback, (\ref{eq: closureMarginal}) becomes exactly
(\ref{eq: closureMarketDensity}), so uniqueness gives $\bar p=p_*$.

For the second feedback use $2n=(q-s^2)p_*$. Subtract the market equation from
(\ref{eq: closureMarginal}) and put $w=\bar p-p_*$. Then
\begin{equation}\label{eq: closureDifference}
\partial_t w=\mathcal B_t w+
       \tfrac12(\partial_{xx}+\partial_x)(s^2w),\qquad w(0)=0.
\end{equation}
Uniqueness for this linear equation gives $w=0$. Hence the first feedback holds
as well. Finally
\[
C_m(t,K)=\int_{\mathbb R}(e^x-K)^+\bar p(t,x)dx=C(t,K).
\]
Conversely, a calibrated coefficient satisfying these assumptions has
$\bar p=p_*$ and satisfies the first feedback by the Malliavin calibration
identity. Its stochastic moments solve (\ref{eq: closureFive}), so they also
satisfy the second feedback.
\end{proof}

The known-denominator version removes division by an \emph{unknown} marginal
from this deterministic formulation. It does not remove small market densities
in the wings. Nor does replacing the marginal by a market density automatically
give an unbiased or consistent finite-particle estimator.

\paragraph{Analytical significance and remaining obstruction.}
The full discounted joint law of $(Z,Y)$ would require seven state coordinates.
For the calibration functional it can be replaced exactly by five fields on
the three physical coordinates. Discounting appears as killing, so an extra
discount-factor coordinate is unnecessary. For a prescribed coefficient,
Proposition \ref{prop: closureIdentification} gives a deterministic and
stochastic identification theorem with nonflat volatility.
For calibration, the known-denominator version is a useful formulation for
a fixed-point or perturbation argument. But it still contains a square root,
an unbounded rate domain, a time-zero Dirac law, and derivatives of the unknown
coefficient. In fact its divergence terms generally differentiate
$g=\partial_xs$ again. Positivity of the variance gap, control of small
densities and nonlinear estimates remain necessary. No contraction or
nonlinear parabolicity conclusion follows from the linear closure alone.

\paragraph{Checks and numerical scope.}
The companion script \texttt{checks/check\_moment\_closure.py} checks the
affine generator identity symbolically, checks both marginal PDE reductions,
and tests weak balances in a prescribed nonflat Monte Carlo model. These are
checks of the representation, not a three-dimensional PDE solver or a smile
recovery experiment. A solver should preserve discounted mass, monitor the
full-state compatibility residuals (\ref{eq: closureConstraints}), verify
the shared-noise fluxes, and compare recovered call prices independently.
